\begin{frame}{Marking \underline{global} inconsistencies}
  Consider this program:
  \[%
    \ELam{f}{\TUnknown}{\EAp{f}{(\EPlus{f}{1})}}
  \]%

  \pause
  $\assignType{f}{\TUnknown}$, so the bidirectional system operates gradually, \pause
  but $f$ is 
  %
  \pause
  \begin{enumerate}
    \item applied as a function \pause (a function?)
      \pause
    \item an operand of $\ECOpPlus$ \pause (a number?)
  \end{enumerate}

  \pause
  \vspace{1em}
  \emph{Global, constraint-based type checking would have caught this!}
\end{frame}

\begin{frame}{Layers upon layers}
  Get the best of both worlds by layering \textbf{constraint-based inference} 
  atop the marked lambda calculus

  \vspace{1em}
  \pause
  \begin{itemize}
    \item The marked lambda calculus localizes and recovers predictably from \emph{local inconsistencies}

      \pause
    \item \textbf{Type hole inference} solves and marks \emph{global inconsistencies}
  \end{itemize}
\end{frame}

\begin{frame}{Type hole inference}
  \begin{itemize}
    \item Occurrences of $\TUnknown$ arise from
      \pause (1) \emph{type holes}, \eg, in $\ELam{x}{\TUnknown}{\cdots}$
      \pause or (2) error marks

      \pause
    \item Gather constraints from the bidirectional systems\pause,
      treating occurrences of $\TUnknown$ as unification variables
      {\footnotesize\parencite{siek2008}}

      \pause
    \item Solve them via a standard unification algorithm {\footnotesize\parencite{huet1976}} \\
      \pause whilst \emph{accumulating the potential types for each variable}

      \pause
    \item Not meant to be a standalone type checker!
  \end{itemize}
\end{frame}

\newcommand{\Id}[1]{\ensuremath{\goodcolor{hole}{#1}}}
\newcommand{\IdU}{\ensuremath{\Id{u}}}
\newcommand{\ProvU}{\ensuremath{\IdU}}
\newcommand{\ProvExp}[1]{\ensuremath{exp\goodcolor{\colorOkSideJudge}{(}#1\goodcolor{\colorOkSideJudge}{)}}}
\newcommand{\ProvIn}[1]{\ensuremath{\to_L\goodcolor{\colorOkSideJudge}{(}#1\goodcolor{\colorOkSideJudge}{)}}}
\newcommand{\ProvOut}[1]{\ensuremath{\to_R\goodcolor{\colorOkSideJudge}{(}#1\goodcolor{\colorOkSideJudge}{)}}}

\begin{frame}{Provenances for types}
  Add type provenances $\Provp$ to the \emph{marked language}:

  \vspace{1em}
  \pause
  \begin{itemize}
    \item Link occurrences of $\TUnknown$ to the originating type hole or error mark
      (denoted by unique id $\IdU$)

    \item Distinguish between different $\TUnknown$ based on locus
  \end{itemize}

  \pause
  \[\begin{array}{rcl}
    \Provp & \Coloneqq & \ProvU \mid \ProvExp{\ProvU} \mid \ProvIn{\Provp} \mid \ProvOut{\Provp} \\
    \pause
    \TMV   & \Coloneqq & \cdots \mid \TUnknown^{p}\\
    \pause
    \ECMV  & \Coloneqq & \cdots \mid \ECFreeId{x}{\IdU} \mid \ECApSynNonMatchedArrowId{\ECMV}{\IdU} \mid \ECInconTypeId{\ECMV}{\IdU}
  \end{array}\]
\end{frame}

\begin{frame}{Generating constraints bidirectionally}
  Augment the type system (of the marked language) to generate sets $C$ of constraints
  $\constrain{\TMV_1}{\TMV_2}$\pause, which force $\TMV_1$ and $\TMV_2$ to be consistent:
  \pause
  \[%
    \synConstraint{\ctx}{\ECMV}{\TMV}{C} \qquad \anaConstraint{\ctx}{\ECMV}{\TMV}{C}
  \]%

  \pause
  \begin{mathpar}
    \judgment{
      \synConstraint{\ctx}{\ECMV}{\TMV'}{C} \\
      \consistent{\TMV}{\TMV'}
    }{
      \anaConstraint{\ctx}{\ECMV}{\TMV}{C \cup \{ \constrain{\TMV}{\TMV'} \}}
    }{MASubsume-C}

    \pause
    \judgment{
        \synConstraint{\ctx}{\ECMV}{\TMV'}{C} \\
        \inconsistent{\TMV}{\TMV'}
    }{
      \anaConstraint{\ctx}{\ECInconTypeId{\ECMV}{u}}{\TMV}{C \cup \{ \constrain{\TMV}{\TUnknown^{\ProvExp{\IdU}}} \}}
    }{MAInconsistentTypes-C}
  \end{mathpar}
\end{frame}

\begin{frame}{Unification with inconsistencies}
  For each occurrence of $\TUnknown$, accumulate the $\mathsf{PotentialTypeSet}$\pause: \\
  \emph{all} potential type fillings as inferred from the constraints

  \[\begin{array}{rrcl}
    \mathsf{PotentialTypeSet} & s & \Coloneqq & \{ t^{*} \} \\
    \mathsf{PotentialType}    & t & \Coloneqq & \TUnknown^\Provp \mid \TNum \mid \TBool \mid \TArrow{s}{s}
  \end{array}\]

  \vspace{1em}
  \pause
  To unify $\constrain{\TMV_1}{\TMV_2}$, recursively merge $\mathsf{PotentialTypeSet(\TMV_1)}$ and
  $\mathsf{PotentialTypeSet(\TMV_2)}$ without substituting
\end{frame}

\newcommand{\PTS}{\operatorname{\mathsf{PotentialTypeSet}}}
\newcommand{\cursor}{\ensuremath{\goodcolor{\colorOkSideJudge}{\bm{\wedge}}}}

\begin{frame}
  \[%
    \ELam{f}{\TUnknown^{\Id{1}}}{\EAp{f}}{(\EPlus{f}{1})}
  \]%

  \uncover<2->{
    yields the (inconsistent) constraints
    \[%
      \{ \onderset<6>{\cursor}{\constrain{\TUnknown^{\Id{1}}}{\TArrow{\TUnknown^{\ProvIn{\Id{1}}}}{\TUnknown^{\ProvOut{\Id{1}}}}}}
       \only<3->{, \onderset<7>{\cursor}{\constrain{\TNum}{\TUnknown^{\Id{1}}}}}
       \only<4->{, \onderset<8>{\cursor}{\constrain{\TNum}{\TUnknown^{\ProvIn{\Id{1}}}}}} \}
    \]%
  }
  
  \vspace{-1em}
  \uncover<5->{
    and solving gives
    \[%
      \PTS(\TUnknown^{\Id{1}})
        =
      \Alt<6->{
        \Alt<7->{
          \Alt<8->{
            \{ \TArrow{\{ \onderline<8>{\TNum} \}}{\{ \TUnknown^{\ProvOut{\Id{1}}} \}},
               \TNum \}
          }{
            \{ \TArrow{\{ \TUnknown^{\ProvIn{\Id{1}}} \}}{\{ \TUnknown^{\ProvOut{\Id{1}}} \}},
             \underline{\TNum} \}
          }
        }{
          \{ \underline{\TArrow{\{ \TUnknown^{\ProvIn{\Id{1}}} \}}{\{ \TUnknown^{\ProvOut{\Id{1}}} \}}} \}
        }
      }{
        \{ \}
      }
    \]%
  }
\end{frame}

\begin{frame}
  \[%
    \ELam{f}{\TUnknown^{\Id{1}}}{\EAp{f}}{(\EPlus{f}{1})}
  \]%
  \[%
    \PTS(\TUnknown^{\Id{1}})
      = \{ \TArrow{\{ \TNum \}}{\{ \TUnknown^{\ProvOut{\Id{1}}} \}},
         \TNum \}
  \]%

  \vspace{1em}
  Now we know:

  \vspace{1em}
  \pause
  \begin{itemize}
    \item Two \emph{potential fillings} of $\TUnknown^{\Id{1}}$ are $\TArrow{\TNum}{\TUnknown}$ and $\TNum$\pause, \\
      which is more than one!

      \pause
    \item $\TUnknown^{\Id{1}}$ is unsolved; inconsistent constraints!
  \end{itemize}
\end{frame}

\begin{frame}[fragile]
  \begin{center}
    \alt<2->{
      \alt<3->{
        \includegraphics[width=0.9\textwidth]{inference-arrow.png}
      }{
        \includegraphics[width=0.9\textwidth]{inference-int.png}
      }
    }{
      \includegraphics[width=0.9\textwidth]{inference.png}
    }
    \vspace{2em}
    \alt<2->{
      \alt<3->{
        \includegraphics[width=0.9\textwidth]{inspector-arrow.png}
      }{
        \includegraphics[width=0.9\textwidth]{inspector-int.png}
      }
    }{
      \includegraphics[width=0.9\textwidth]{inspector.png}
    }
  \end{center}

  Hazel offers the user both choices and the ability to fill either in

  \begin{tikzpicture}[remember picture,overlay]
    \node[xshift=-4.6cm, yshift=-5cm, visible on=<2>] at (current page.north east) { \mousePointer };
    \node[xshift=-2.3cm, yshift=-5cm, visible on=<3>] at (current page.north east) { \mousePointer };
  \end{tikzpicture}
\end{frame}

\begin{frame}
  This approach is \emph{neutral}

  \vspace{1em}
  \pause
  \begin{itemize}
    \item Localize errors to the originating type hole or error mark\pause, \\
      instead of guessing about user intent

      \pause
    \item All potential type hole fillings are provided to the user for selection

      \pause
    \item Control returned to bidirectional system after user selects
  \end{itemize}
\end{frame}
